\begin{landscape}
\begingroup
\fontsize{10}{13}\selectfont
\fancyhead[L]{\small NACHLASS. STÖRUNGEN DER PALLAS}
\renewcommand{\headrulewidth}{0pt}
\par\medskip
\pdfbookmark[2]{Nachlass VI [5]}{pallas-VI-fragment5}
\begin{center}[5.]\end{center}
\par\noindent\hspace*{\parindent}[Die Störungen der großen Halbachse ergeben sich genau auf dieselbe Weise wie in der ersten Rechnung. Für unser Beispiel erhält man die folgende Rechnung, wenn man die Konstante \(=-\tfrac{2}{8}\lambda\dfrac{10000000}{206265}\) setzt, also die 7.\ Dezimalstelle als Einheit wählt:
\SourceLeaf{556}
\par
\begingroup\small
\noindent\begin{minipage}[c]{.43\textwidth}
\begin{tabular}{@{}lr@{}}
\(b\)\enspace .\,.\,.\,.\,.\,.\,.\,.\,.\,.\,.\,.\,.&\(0{,}95337\)\\
\(b'\)\enspace .\,.\,.\,.\,.\,.\,.\,.\,.\,.\,.\,.\,.&\(0{,}52101\)\\
\multicolumn{2}{r}{\(Q=69^0\,43'\,11''\)}\\
\(\sin Q\)\enspace .\,.\,.\,.\,.\,.\,.\,.\,.&\(9{,}97221\)\\
\(q\)\enspace .\,.\,.\,.\,.\,.\,.\,.\,.\,.\,.\,.\,.&\(0{,}98116\)\\
\(2-5\dfrac{n'}{n}\)\enspace .\,.\,.\,.\,.\,.&\(8{,}74430\)\\
\hline
\(*\)\enspace .\,.\,.\,.\,.\,.\,.\,.\,.\,.\,.\,.\,.&\(2{,}23686\)\\
Konstante\enspace .\,.\,.\,.\,.\,.\,.&\(1{,}14727_n\)\\
\hline
&\(3{,}38413_n\)
\end{tabular}
\end{minipage}\hfill
\begin{minipage}[c]{.55\textwidth}
also das Störungsglied
\[
-2422\cos(5M'-2M+69^0\,43'\,11'').
\]
\end{minipage}\par
\endgroup
\par\noindent\hspace*{\parindent}Das vollständige Ergebnis ist in der folgenden Notiz zusammen mit den Störungen des ersten Teils der Epoche angegeben, da die Argumente dort dieselben sind. Dabei ist stets der Sinus des betreffenden Arguments zu nehmen, sodass z.\ B. das eben berechnete Störungsglied lautet:
\[
+2422\sin(5\text{\mdseries\AstroSymbols ♃}-2\text{\mdseries\AstroSymbols ⚴}+165^0\,35'\,33'').]
\]
\clearpage
\endgroup
\end{landscape}
